\documentclass[11pt,a4paper]{article}
\input{preamble}
\title{Stochastic Financial Models}
\begin{document}
\maketitle
\thispagestyle{empty}
\section{Martingales}

We work with a probability space $(\Omega,\mathcal{F},\mathbb{P})$.

\subsection{Definitions}

A \emph{filtration} is a sequence of $\sigma$-algebras
$(\mathcal{F}_n)_{n\geq0}$ such that
\[
\mathcal{F}_n\subseteq\mathcal{F}_{n+1}\subseteq\mathcal{F}.
\]

A \emph{random process} is a sequence
$X=(X_n)_{n\geq0}$ of random variables.

Given a filtration $(\mathcal{F}_n)_{n\geq0}$ on
$(\Omega,\mathcal{F},\mathbb{P})$, we say that $X$ is a
\emph{martingale} if:
\begin{enumerate}
    \item[(a)] $X$ is adapted, i.e.\ $X_n$ is
    $\mathcal{F}_n$-measurable for all $n$.

    \item[(b)] $X$ is integrable, i.e.\ $X_n$ is integrable
    for all $n$:
    \[
    \mathbb{E}[|X_n|]<\infty.
    \]

    \item[(c)] $X$ satisfies the martingale property:
    \[
    \mathbb{E}[X_{n+1}\mid\mathcal{F}_n]
    =
    X_n
    \qquad\text{a.s., for all }n.
    \]
\end{enumerate}

Given (a) and (b), condition (c) is equivalent to
\[
\mathbb{E}[X_{n+1}\mathbf{1}_A]
=
\mathbb{E}[X_n\mathbf{1}_A]
\qquad
\text{for all }n
\text{ and all }A\in\mathcal{F}_n.
\]

If (c) is replaced by
\[
\mathbb{E}[X_{n+1}\mid\mathcal{F}_n]
\geq X_n
\qquad\text{a.s., for all }n,
\]
or, equivalently,
\[
\mathbb{E}[X_{n+1}\mathbf{1}_A]
\geq
\mathbb{E}[X_n\mathbf{1}_A]
\qquad
\text{for all }n
\text{ and all }A\in\mathcal{F}_n,
\]
then $(X_n)_{n\geq0}$ is a \emph{submartingale}.

If (c) is replaced by
\[
\mathbb{E}[X_{n+1}\mid\mathcal{F}_n]
\leq X_n
\qquad\text{a.s., for all }n,
\]
or, equivalently,
\[
\mathbb{E}[X_{n+1}\mathbf{1}_A]
\leq
\mathbb{E}[X_n\mathbf{1}_A]
\qquad
\text{for all }n
\text{ and all }A\in\mathcal{F}_n,
\]
then $(X_n)_{n\geq0}$ is a \emph{supermartingale}.

\subsection{Natural filtration}

Given a random process $X=(X_n)_{n\geq0}$, its
\emph{natural filtration} $(\mathcal{F}_n^X)_{n\geq0}$
is defined by
\[
\begin{aligned}
\mathcal{F}_n^X
&=
\sigma(X_0,X_1,\ldots,X_n)\\
&=
\left\{
\{(X_0,X_1,\ldots,X_n)\in B\}:
B\in\mathcal{B}(\mathbb{R}^{n+1})
\right\}.
\end{aligned}
\]
This is the smallest filtration making $X$ adapted.

\paragraph{Exercise.}
Show that a martingale $X$ with respect to
$(\mathcal{F}_n)_{n\geq0}$ is also a martingale with respect
to $(\mathcal{F}_n^X)_{n\geq0}$.

\subsection{Examples}

Let $X_1,X_2,\ldots$ be a sequence of independent integrable
random variables. Set
\[
\mathcal{F}_n=\mathcal{F}_n^X.
\]
Define
\[
S_0=0,
\qquad
Z_0=1,
\]
and
\[
S_n=X_1+\cdots+X_n,
\qquad
Z_n=\prod_{i=1}^{n}X_i.
\]

\begin{itemize}
    \item If $\mathbb{E}[X_n]=0$ for all $n$, then
    $(S_n)_{n\geq0}$ is a martingale.

    \item If $X_n\geq0$ and $\mathbb{E}[X_n]=1$ for all $n$,
    then $(Z_n)_{n\geq0}$ is a martingale.
\end{itemize}

To see this, take $A\in\mathcal{F}_n$. Then
\[
\begin{aligned}
\mathbb{E}[S_{n+1}\mathbf{1}_A]
&=
\mathbb{E}[X_{n+1}\mathbf{1}_A]
+
\mathbb{E}[S_n\mathbf{1}_A]\\
&=
\underbrace{
\mathbb{E}[X_{n+1}]\mathbb{P}(A)
}_{=0}
+
\mathbb{E}[S_n\mathbf{1}_A]\\
&=
\mathbb{E}[S_n\mathbf{1}_A].
\end{aligned}
\]

Also,
\[
\begin{aligned}
\mathbb{E}[Z_{n+1}\mathbf{1}_A]
&=
\mathbb{E}[X_{n+1}Z_n\mathbf{1}_A]\\
&=
\underbrace{
\mathbb{E}[X_{n+1}]
}_{=1}
\mathbb{E}[Z_n\mathbf{1}_A]\\
&=
\mathbb{E}[Z_n\mathbf{1}_A].
\end{aligned}
\]
This also shows integrability.

\section{Conditional Probabilities and Expectation}

For $\mathbb{P}(B)>0$,
\[
\mathbb{P}(A\mid B)
=
\frac{\mathbb{P}(A\cap B)}{\mathbb{P}(B)},
\]
and
\[
\mathbb{E}[X\mid B]
=
\frac{\mathbb{E}[X\mathbf{1}_B]}{\mathbb{P}(B)}.
\]

\subsection{Existence and uniqueness}

\paragraph{Theorem.}
Let $\mathcal{G}\subseteq\mathcal{F}$ be a $\sigma$-algebra.
Let $X$ be an integrable random variable. Then there exists
a random variable $Y$ such that:
\begin{enumerate}
    \item[(a)] $Y$ is $\mathcal{G}$-measurable;

    \item[(b)]
    \[
    \mathbb{E}[X\mathbf{1}_A]
    =
    \mathbb{E}[Y\mathbf{1}_A]
    \qquad\text{for all }A\in\mathcal{G}.
    \]
\end{enumerate}
Moreover, if $Y'$ is another random variable satisfying
(a) and (b), then
\[
Y=Y'\qquad\text{a.s.}
\]

\paragraph{Proof of uniqueness.}
Let $X'$ be another integrable random variable with
$X\leq X'$. Suppose $Y'$ satisfies (a) and (b) with respect
to $X'$.

Set
\[
A=\{Y>Y'\},
\qquad
Z=(Y-Y')\mathbf{1}_A.
\]
Then $A\in\mathcal{G}$ and $Z\geq0$. Hence
\[
\begin{aligned}
\mathbb{E}[Z]
&=
\mathbb{E}[(Y-Y')\mathbf{1}_A]\\
&=
\mathbb{E}[Y\mathbf{1}_A]
-
\mathbb{E}[Y'\mathbf{1}_A]\\
&=
\mathbb{E}[X\mathbf{1}_A]
-
\mathbb{E}[X'\mathbf{1}_A]\\
&=
\mathbb{E}[(X-X')\mathbf{1}_A]
\leq0.
\end{aligned}
\]
Therefore $Z=0$ a.s., so
\[
Y\leq Y'\qquad\text{a.s.}
\]
Taking $X'=X$, we also obtain $Y\geq Y'$ a.s. Thus
\[
Y=Y'\qquad\text{a.s.}
\]

\subsection{Notation}

We write
\[
Y=\mathbb{E}[X\mid\mathcal{G}]
\qquad\text{a.s.}
\]

If $\mathcal{G}=\sigma(W)$, where $W$ is a random variable,
we write
\[
\mathbb{E}[X\mid\mathcal{G}]
=
\mathbb{E}[X\mid W].
\]

If $X=\mathbf{1}_A$, we write
\[
\mathbb{E}[X\mid\mathcal{G}]
=
\mathbb{P}(A\mid\mathcal{G}).
\]
\\

\begin{center}
    \textbf{End of Lecture 1}
\end{center}

\section*{Optional stopping}

We say that a random time
\[
T:\Omega\longrightarrow \{0,1,2,\ldots\}\cup\{\infty\}
\]
is a \textbf{stopping time} if
\[
\{T\leq n\}\in\mathcal{F}_n
\qquad\text{for all }n.
\]

\noindent\textbf{Example.}
Let $(X_n)_{n\geq 0}$ be an adapted process, with values in $\mathbb{Z}$, say.
Set
\[
T=\inf\{n\geq 0:X_n=0\}.
\]
Then
\[
\{T\leq n\}
=
\underbrace{\{X_0=0\}}_{\in\mathcal{F}_0}
\cup
\underbrace{\{X_1=0\}}_{\in\mathcal{F}_1}
\cup\cdots\cup
\underbrace{\{X_n=0\}}_{\in\mathcal{F}_n},
\]
with
\[
\mathcal{F}_0\subseteq\mathcal{F}_1
\subseteq\cdots\subseteq\mathcal{F}_n.
\]
So $T$ is a stopping time.

\subsection*{Optional stopping theorem}

Let $(M_n)_{n\geq 0}$ be a martingale and let $T$ be a bounded stopping time.
Then
\[
\mathbb{E}(M_T)=\mathbb{E}(M_0),
\qquad
\bigl(M_T(\omega)=M_{T(\omega)}(\omega)\bigr).
\]

\noindent\textbf{Proof.}
Choose $n$ such that $T\leq n$. Then
\[
\begin{aligned}
M_T
&=M_0+(M_1-M_0)+\cdots+(M_T-M_{T-1})\\
&=M_0+\sum_{k=1}^{n}(M_k-M_{k-1})
  \mathbf{1}_{\{T\geq k\}}.
\end{aligned}
\]
Now
\[
\{T\geq k\}
=
\{T\leq k-1\}^{c}
\in\mathcal{F}_{k-1},
\]
so
\[
\mathbb{E}\!\left(
(M_k-M_{k-1})\mathbf{1}_{\{T\geq k\}}
\right)=0
\qquad\text{for all }k.
\]
Take expectations to see
\[
\mathbb{E}(M_T)=\mathbb{E}(M_0)+0.
\qquad\square
\]

\subsection*{OST II}

Let $(M_n)_{n\geq 0}$ be a martingale and let $T$ be a finite-valued stopping time.
Suppose there exists $C<\infty$ such that one of the following holds:
\begin{enumerate}
    \item[(a)]
    \[
    |M_n|\leq C
    \qquad\text{for all }n\leq T.
    \]

    \item[(b)]
    \[
    \mathbb{E}(T)<\infty
    \qquad\text{and}\qquad
    |M_n-M_{n-1}|\leq C
    \qquad\text{for all }n\leq T.
    \]
\end{enumerate}
Then
\[
\mathbb{E}(M_T)=\mathbb{E}(M_0).
\]

\noindent\textbf{Proof.}
Define
\[
T_n=T\wedge n.
\]
Then $T_n$ is a bounded stopping time.
Since $T<\infty$, $T_n=T$ eventually, so
\[
M_{T_n}\longrightarrow M_T
\qquad\text{as }n\longrightarrow\infty.
\]
By optional stopping, we have
\[
\mathbb{E}(M_{T_n})=\mathbb{E}(M_0).
\]

If (a) holds, then
\[
\sup_n |M_{T_n}|\leq C,
\]
so
\[
\mathbb{E}(M_{T_n})\longrightarrow\mathbb{E}(M_T)
\]
by bounded convergence.

If (b) holds, then
\[
\sup_n |M_{T_n}|\leq CT+|M_0|
\qquad\text{and}\qquad
\mathbb{E}(T)<\infty,
\]
so
\[
\mathbb{E}(M_{T_n})\longrightarrow\mathbb{E}(M_T)
\]
by dominated convergence.

In any case, we can let $n\longrightarrow\infty$ to see
\[
\mathbb{E}(M_T)=\mathbb{E}(M_0).
\qquad\square
\]

\subsection*{Examples}

Let $X_1,X_2,\ldots$ and $Y_1,Y_2,\ldots$ be independent random variables such that
\[
\mathbb{P}(X_n=1)=\mathbb{P}(X_n=-1)=\frac12,
\qquad
\mathbb{P}(Y_n=0)=\mathbb{P}(Y_n=2)=\frac12.
\]
Set
\[
S_n=1+X_1+\cdots+X_n,
\qquad
Z_n=\prod_{i=1}^{n}Y_i,
\qquad
(Z_0=1).
\]
Then $(S_n)_{n\geq 0}$ and $(Z_n)_{n\geq 0}$ are martingales.

Set
\[
T=\inf\{n\geq 0:S_n=0\}.
\]
Then
\[
\mathbb{P}(T<\infty)=1,
\]
but
\[
\mathbb{E}(S_0)=1,
\qquad
\mathbb{E}(S_T)=0.
\]

Now set
\[
T=\inf\{n\geq 0:Z_n=0\}.
\]
Then
\[
\mathbb{P}(T\leq n)=1-2^{-n},
\]
so
\[
\mathbb{P}(T<\infty)=1.
\]
In fact,
\[
\mathbb{E}(T)<\infty.
\]
But
\[
\mathbb{E}(Z_T)=0,
\qquad
\mathbb{E}(Z_0)=1.
\]

\medskip

Recall that, for a $\sigma$-algebra $\mathcal{G}\subseteq\mathcal{F}$ and an integrable random variable $X$,
$\mathbb{E}(X\mid\mathcal{G})$ is characterized by
\[
Y=\mathbb{E}(X\mid\mathcal{G})
\qquad\text{a.s.}
\]
if
\begin{enumerate}
    \item[(a)] $Y$ is $\mathcal{G}$-measurable;

    \item[(b)] $Y$ is integrable and
    \[
    \mathbb{E}(X\mathbf{1}_A)
    =
    \mathbb{E}(Y\mathbf{1}_A)
    \qquad\text{for all }A\in\mathcal{G}.
    \]
\end{enumerate}

\section*{Properties of conditional expectation}

\begin{enumerate}
    \item[(i)]
    \[
    \mathbb{E}\!\left(\mathbb{E}(X\mid\mathcal{G})\right)
    =
    \mathbb{E}(X).
    \]

    \item[(ii)]
    If $X$ is $\mathcal{G}$-measurable, then
    \[
    \mathbb{E}(X\mid\mathcal{G})=X
    \qquad\text{a.s.}
    \]

    \item[(iii)]
    If $X$ is independent of $\mathcal{G}$, then
    \[
    \mathbb{E}(X\mid\mathcal{G})=\mathbb{E}(X)
    \qquad\text{a.s.}
    \]

    \item[(iv)]
    If $X\geq 0$ a.s., then
    \[
    \mathbb{E}(X\mid\mathcal{G})\geq 0
    \qquad\text{a.s.}
    \]

    \item[(v)]
    \textbf{Linearity.}
    For $W,X$ integrable and $\alpha,\beta\in\mathbb{R}$,
    \[
    \mathbb{E}(\alpha X+\beta W\mid\mathcal{G})
    =
    \alpha\mathbb{E}(X\mid\mathcal{G})
    +
    \beta\mathbb{E}(W\mid\mathcal{G}).
    \]
\end{enumerate}

Take
\[
Y=\mathbb{E}(X\mid\mathcal{G}),
\qquad
U=\mathbb{E}(W\mid\mathcal{G}).
\]
Then $\alpha Y+\beta U$ is $\mathcal{G}$-measurable and integrable.

For $A\in\mathcal{G}$,
\[
\begin{aligned}
\mathbb{E}\!\left((\alpha X+\beta W)\mathbf{1}_A\right)
&=\alpha\mathbb{E}(X\mathbf{1}_A)
  +\beta\mathbb{E}(W\mathbf{1}_A)\\
&=\alpha\mathbb{E}(Y\mathbf{1}_A)
  +\beta\mathbb{E}(U\mathbf{1}_A)\\
&=\mathbb{E}\!\left((\alpha Y+\beta U)\mathbf{1}_A\right).
\end{aligned}
\]
Hence
\[
\alpha Y+\beta U
=
\mathbb{E}(\alpha X+\beta W\mid\mathcal{G})
\qquad\text{a.s.}
\]

\subsection*{Proposition (Tower property)}

Let $\mathcal{H}$ be another $\sigma$-algebra, with
\[
\mathcal{H}\subseteq\mathcal{G}.
\]
Then
\[
\mathbb{E}\!\left(
\mathbb{E}(X\mid\mathcal{G})\mid\mathcal{H}
\right)
=
\mathbb{E}(X\mid\mathcal{H}).
\]

\noindent\textbf{Proof.}
Take
\[
Y=
\mathbb{E}\!\left(
\mathbb{E}(X\mid\mathcal{G})\mid\mathcal{H}
\right)
\qquad\text{a.s.}
\]
Then, for $A\in\mathcal{H}$, we have $A\in\mathcal{G}$, so
\[
\mathbb{E}(Y\mathbf{1}_A)
\underset{A\in\mathcal{H}}{=}
\mathbb{E}\!\left(
\mathbb{E}(X\mid\mathcal{G})\mathbf{1}_A
\right)
\underset{A\in\mathcal{G}}{=}
\mathbb{E}(X\mathbf{1}_A).
\]
$Y$ is integrable and $\mathcal{H}$-measurable, so
\[
Y=\mathbb{E}(X\mid\mathcal{H})
\qquad\text{a.s.}
\qquad\square
\]

\begin{center}
\textbf{End of Lecture 2}
\end{center}
\end{document}
